\begin{tikzpicture}[
  x=1mm,y=1mm,
  >=Stealth,
  font=\figurefont\fontsize{8}{10}\selectfont
]

\tikzset{
  vnode/.style={circle,draw=FigureBlue,fill=FigureBlueFill,minimum size=4.4mm,inner sep=0pt,line width=.55pt,text=FigureInk},
  edge/.style={draw=FigureStroke,line width=.6pt},
  cutedge/.style={draw=FigureRed,line width=.85pt},
  cutline/.style={dashed,draw=FigureRed,line width=.65pt},
  titlelbl/.style={font=\figurefont\fontsize{8}{10}\selectfont,text=FigureInk},
  sublbl/.style={font=\figurefont\fontsize{7.5}{9}\selectfont,text=FigureMuted}
}

% Top-left: (a) Path graph P_5
\begin{scope}[xshift=0mm,yshift=28mm]
  \node[vnode] (p1) at (0,8) {};
  \node[vnode] (p2) at (8,8) {};
  \node[vnode] (p3) at (16,8) {};
  \node[vnode] (p4) at (24,8) {};
  \node[vnode] (p5) at (32,8) {};
  \draw[edge] (p1)--(p2);
  \draw[cutedge] (p2)--(p3);
  \draw[edge] (p3)--(p4)--(p5);
  \draw[cutline] (12,1.5) -- (12,14.5);
  \node[titlelbl] at (16,-2.5) {(a) 链状图 $P_5$};
  \node[sublbl] at (16,-7.2) {单边即断，跨割边数 $1$};
\end{scope}

% Top-right: (b) Barbell graph (two K_4 connected by a bridge)
\begin{scope}[xshift=54mm,yshift=28mm]
  \node[vnode] (b1) at (0,13) {};
  \node[vnode] (b2) at (0,3) {};
  \node[vnode] (b3) at (9,13) {};
  \node[vnode] (b4) at (9,3) {};
  \node[vnode] (b5) at (21,13) {};
  \node[vnode] (b6) at (21,3) {};
  \node[vnode] (b7) at (30,13) {};
  \node[vnode] (b8) at (30,3) {};
  % Left K_4
  \draw[edge] (b1)--(b2)--(b4)--(b3)--(b1);
  \draw[edge] (b1)--(b4);
  \draw[edge] (b2)--(b3);
  % Right K_4
  \draw[edge] (b5)--(b6)--(b8)--(b7)--(b5);
  \draw[edge] (b5)--(b8);
  \draw[edge] (b6)--(b7);
  % Bridge
  \draw[cutedge] (b3)--(b5);
  \draw[cutline] (15,0.5) -- (15,15.5);
  \node[titlelbl] at (15,-2.5) {(b) 杠铃图（单桥连接）};
  \node[sublbl] at (15,-7.2) {团内稠密，桥连接两侧，跨割边数 $1$};
\end{scope}

% Bottom-left: (c) Grid graph 3x3
\begin{scope}[xshift=4mm,yshift=0mm]
  \foreach \ix in {0,1,2}{
    \foreach \iy in {0,1,2}{
      \node[vnode] (g\ix\iy) at (\ix*11,\iy*6.5+1.5) {};
    }
  }
  \foreach \iy in {0,1,2}{
    \draw[cutedge] (g0\iy)--(g1\iy);
    \draw[edge] (g1\iy)--(g2\iy);
  }
  \foreach \ix in {0,1,2}{
    \draw[edge] (g\ix0)--(g\ix1)--(g\ix2);
  }
  \draw[cutline] (5.5,-0.5) -- (5.5,16.5);
  \node[titlelbl] at (12,-2.5) {(c) 网格图 $P_3\mathbin{\mdlgwhtsquare} P_3$（多路冗余）};
  \node[sublbl] at (12,-7.2) {切口横跨整列，跨割边数 $3$};
\end{scope}

% Bottom-right: (d) Complete graph K_5
\begin{scope}[xshift=56mm,yshift=0mm]
  \node[vnode] (k1) at (4,13) {};
  \node[vnode] (k2) at (4,3) {};
  \node[vnode] (k3) at (20,15) {};
  \node[vnode] (k4) at (25,8) {};
  \node[vnode] (k5) at (20,1) {};
  % Internal edges
  \draw[edge] (k1)--(k2);
  \draw[edge] (k3)--(k4)--(k5)--(k3);
  % Cross edges (2 x 3 = 6 edges)
  \foreach \u in {k1,k2}{
    \foreach \v in {k3,k4,k5}{
      \draw[cutedge] (\u)--(\v);
    }
  }
  \draw[cutline] (11.5,-0.5) -- (11.5,16.5);
  \node[titlelbl] at (13,-2.5) {(d) 完全图 $K_5$};
  \node[sublbl] at (13,-7.2) {任意两点直连，跨割边数 $6$};
\end{scope}
\end{tikzpicture}
